\chapter{Saludos y presentaciones}
\section{Objetivos}
\begin{itemize}[leftmargin=*]
	\item Usar saludos formales e informales y despedidas.
	\item Presentarte y pedir información básica (nombre, procedencia, idioma).
	\item Practicar cortesía: \textit{alstublieft/alsjeblieft}, \textit{dank u/je}.
\end{itemize}

\section{Contenidos clave}
\begin{itemize}[leftmargin=*]
	\item Saludos y despedidas formales vs. informales.
	\item Presentación básica: nombre, origen, idiomas.
	\item Cortesía mínima en tienda y primer contacto.
	\item Pronunciación y entonación en intercambios cortos.
\end{itemize}

\section{Ruta del capítulo}
Ocho secciones breves con ejemplos y mini práctica. Úsalas en orden.

\section{Saludos formales e informales}
\textbf{Qué son:} Fórmulas para abrir interacción según contexto.

\textbf{Por qué importan:} Ajustar registro muestra cortesía y evita sonar brusco.

\textbf{Frases clave}
\begin{itemize}[leftmargin=*]
  \item Informal: \textit{Hoi! Hallo! Hey!}
  \item Neutro/formal: \textit{Goeiedag}, \textit{Goedemorgen}, \textit{Goedemiddag}, \textit{Goedenavond}.
\end{itemize}

\textbf{Práctica rápida}
\begin{itemize}[leftmargin=*]
  \item Di 5 saludos al espejo cambiando tono (amistoso vs.\ neutro).
  \item Role-play: llegada a oficina (formal) y a casa de amigos (informal).
\end{itemize}

\section{Despedidas}
\textbf{Qué son:} Fórmulas para cerrar la interacción con amabilidad.

\textbf{Frases clave}
\begin{itemize}[leftmargin=*]
  \item Informal: \textit{Doei! Dag! Tot straks!}
  \item Neutro/formal: \textit{Tot ziens! Fijne dag! Prettige avond!}
\end{itemize}

\textbf{Cómo practicarlas}
\begin{itemize}[leftmargin=*]
  \item Empareja saludo + despedida coherente (ej.: \textit{Goedemorgen} \textrightarrow{} \textit{Tot ziens}).
  \item Graba diálogos de 4 líneas cerrando con despedida adecuada.
\end{itemize}

\section{Presentarse (nombre y origen)}
\textbf{Objetivo:} Decir quién eres y de dónde vienes.

\textbf{Estructuras}
\begin{itemize}[leftmargin=*]
  \item Nombre: \textit{Ik ben	extellipsis{}} / \textit{Ik heet	extellipsis{}}.
  \item Origen: \textit{Ik kom uit	extellipsis{}} / \textit{Ik woon in	extellipsis{}}.
  \item Idioma: \textit{Ik spreek Spaans en een beetje Nederlands}.
\end{itemize}

\textbf{Práctica}
\begin{itemize}[leftmargin=*]
  \item Completa tarjetas con tus datos y preséntate en voz alta.
  \item Cambia un dato y repite (otra ciudad, otro idioma aprendido).
\end{itemize}

\section{Preguntar y decir cómo estás}
\textbf{Objetivo:} Manejar la pregunta de cortesía y respuestas breves.

\textbf{Frases clave}
\begin{itemize}[leftmargin=*]
  \item Pregunta: \textit{Hoe gaat het (met je/u)?}
  \item Respuestas: \textit{Goed, dank je. En met jou?} / \textit{Het gaat wel.} / \textit{Niet zo goed.}
\end{itemize}

\textbf{Práctica}
\begin{itemize}[leftmargin=*]
  \item Intercambia 3 rondas de pregunta-respuesta cambiando el estado.
  \item Añade un motivo: \textit{Ik ben moe, want ik heb veel gewerkt}.
\end{itemize}

\section{Frases de cortesía básicas}
\textbf{Qué son:} Expresiones para pedir, agradecer y disculparse.

\textbf{Frases clave}
\begin{itemize}[leftmargin=*]
  \item \textit{Alsjeblieft/Alstublieft} (informal/formal) para ofrecer o pedir.
  \item \textit{Dank je/Dank u} para agradecer; \textit{Graag gedaan} para responder.
  \item \textit{Sorry, het spijt me} para disculparse.
\end{itemize}

\textbf{Práctica}
\begin{itemize}[leftmargin=*]
  \item Role-play: pedir café (informal) y pedir ayuda en oficina (formal).
  \item Cambia registro: reescribe una petición informal a formal y viceversa.
\end{itemize}

\section{Presentar a otra persona}
\textbf{Objetivo:} Introducir a alguien indicando nombre y relación.

\textbf{Estructuras}
\begin{itemize}[leftmargin=*]
	\item \textit{Dit is\textellipsis{}} (informal), \textit{Mag ik u voorstellen aan\textellipsis{}} (formal).
	\item Añadir relación: \textit{mijn collega, vriend of vriendin, partner}.
\end{itemize}

\textbf{Práctica}
\begin{itemize}[leftmargin=*]
	\item Mini escena: presentas a un colega a tu jefe (formal) y a un amigo (informal).
	\item Añade un dato extra: ciudad u ocupación.
\end{itemize}

% chktex-file 13
\section{Diálogos simples de presentación}
\textbf{Objetivo:} Combinar saludo, nombre, origen y cortesía en un intercambio breve.

\textbf{Modelo guiado}
\begin{quote}
A: \textit{Hoi, ik ben Marta, en jij?}\\
B: \textit{Ik heet Tom, ik kom uit Gent. Waar kom jij vandaan?}\\
A: \textit{Ik kom uit Lima, leuk je te ontmoeten!}
\end{quote}

\textbf{Tareas}
\begin{itemize}[leftmargin=*]
  \item Reescribe el diálogo en registro formal (\textit{goedemiddag, ik heet	extellipsis{}, u}).
  \item Graba dos variaciones cambiando país e idioma.
\end{itemize}

\section{Pronunciación y entonación básica}
\textbf{Foco:} Ritmo descendente en enunciados; subida ligera en preguntas de sí/no.

\textbf{Tips rápidos}
\begin{itemize}[leftmargin=*]
  \item Marca acento en la primera sílaba: \textit{GOEde-mor-gen}, \textit{HAL-lo}, \textit{LEUk}.
  \item Preguntas de cortesía suben al final: \textit{Hoe gaat het?} (subida suave).
  \item Suaviza la \textit{g/ch} en \textit{goed}, \textit{dag} practicando con aire.
\end{itemize}

\textbf{Práctica}
\begin{itemize}[leftmargin=*]
  \item Lee en voz alta 5 saludos y 3 preguntas grabándote; revisa la curva de entonación.
  \item Shadowing de un diálogo corto de presentación (30 s).
\end{itemize}


\section{Vocabulario clave}
\begin{itemize}[leftmargin=*]
	\item Saludos: \textit{goedemorgen}, \textit{goedemiddag}, \textit{goedenavond}, \textit{hallo}, \textit{hoi}.
	\item Despedidas: \textit{dag}, \textit{tot ziens}, \textit{doei}.
	\item Cortesía: \textit{alstublieft/alsjeblieft}, \textit{dank u/je}, \textit{sorry}, \textit{pardon}.
\end{itemize}

\section{Pronunciación}
\begin{itemize}[leftmargin=*]
	\item Alarga las vocales dobles en \textit{goedemorgen} y marca la \textit{g} velar suave.
	\item Entonación descendente en despedidas cortas; ascendente en preguntas de cortesía.
\end{itemize}

\section{Errores comunes}
\begin{itemize}[leftmargin=*]
	\item Confundir formal (\textit{u}) con informal (\textit{je}); ajusta al contexto.
	\item Omitir \textit{alsjeblieft/alstublieft} al pedir algo; añádelo tras la petición.
	\item Usar entonación plana al preguntar; sube el tono final en preguntas cerradas.
\end{itemize}

\section{Práctica guiada}
\begin{itemize}[leftmargin=*]
	\item Lee y graba tres diálogos modelo, cambiando formalidad.
	\item Rellena huecos con saludo/despedida adecuado según la hora.
	\item Intercambio simulado: pides un café y das las gracias.
\end{itemize}

\section{Ejercicios de producción}
\begin{itemize}[leftmargin=*]
	\item Escribe y graba un diálogo de 8 turnos presentándote y preguntando origen/idioma.
	\item Cambia el diálogo a versión formal e informal manteniendo cortesía.
\end{itemize}

\section{Mini repaso}
\begin{itemize}[leftmargin=*]
	\item Distingues cuándo usar \textit{u} vs. \textit{je}.
	\item Produces saludos y despedidas adecuados según la hora.
	\item Mantienes entonación correcta en preguntas de cortesía.
\end{itemize}

\section{Resultado esperado}
\begin{itemize}[leftmargin=*]
	\item Al terminar puedes iniciar y cerrar conversaciones básicas con saludos apropiados, presentarte y pedir datos mínimos con cortesía y buena entonación.\end{itemize}

\section{Microevaluación}
\begin{itemize}[leftmargin=*]
	\item Presentación oral de 6 frases y mini diálogo grabado.
\end{itemize}
